\chapter{Cơ sở lý thuyết (Theoretical Background)}

\section{Information Retrieval Pipeline}
Với hệ thống truy xuất thông tin (Information Retrieval - IR), thay vì thực hiện thao tác rà soát tuần tự trên toàn bộ cơ sở dữ liệu vốn không khả thi với dữ liệu lớn, hệ thống IR hoạt động dựa trên một quy trình xử lý luồng (pipeline) khép kín. Quy trình này chịu trách nhiệm chuyển đổi khối lượng lớn dữ liệu phi cấu trúc thành các ma trận toán học và cấu trúc lưu trữ tối ưu, bao gồm ba giai đoạn chính:
\begin{enumerate}
    \item \textbf{Ingestion (Thu thập và làm sạch):} Dữ liệu thô từ nhiều nguồn được thu thập, tiến hành chuẩn hóa định dạng và loại bỏ các thành phần nhiễu giúp bóc tách và trích xuất các trường thông tin quan trọng nhằm tạo ra một tập dữ liệu đầu vào sạch, đáp ứng yêu cầu dữ liệu đầu vào cho các bước xử lý tiếp theo.
    \item \textbf{Indexing (Lập chỉ mục):} Dữ liệu sau khi làm sạch được tổ chức thành các cấu trúc dữ liệu đặc biệt (như Inverted Index đối với văn bản) để tối ưu hóa tốc độ tìm kiếm.
    \item \textbf{Query Processing (Xử lý truy vấn):} Hệ thống tiếp nhận câu truy vấn từ người dùng, phân tích (tokenize), và đối chiếu với chỉ mục để trả về các kết quả phù hợp nhất dựa trên thuật toán xếp hạng (ranking algorithm).
\end{enumerate}

\section{Inverted Index và Thuật toán BM25}
Trong giai đoạn hiện tại, nhằm đáp ứng yêu cầu tìm kiếm với hiệu năng cao trên các tập dữ liệu lớn, hệ thống sử dụng các công nghệ tìm kiếm văn bản cốt lõi:
\begin{itemize}
    \item \textbf{Inverted Index (Chỉ mục đảo ngược):} Đây là cấu trúc dữ liệu nền tảng của Elasticsearch. Thay vì quét từng tài liệu để tìm từ khóa, Inverted Index lưu trữ một danh sách các từ khóa (terms) và ánh xạ mỗi từ khóa đến danh sách các tài liệu (documents) chứa nó. Điều này giúp tốc độ tìm kiếm đạt mức gần thời gian thực (near real-time) ngay cả với Big Data.
    \item \textbf{BM25 (Best Matching 25):} Là thuật toán xếp hạng mặc định của Elasticsearch, được phát triển và cải tiến dựa trên nền tảng của mô hình trọng số TF-IDF (Term Frequency - Inverse Document Frequency) \cite{robertson2009bm25}. Thuật toán này giúp khai phá văn bản và truy xuất thông tin, đánh giá và xếp hạng độ liên quan của các tài liệu trả về so với câu truy vấn của người dùng dựa trên:
    \begin{itemize}
        \item \textbf{TF (Term Frequency):} Tần suất xuất hiện của từ khóa trong tài liệu. Tần suất càng cao, điểm số càng tăng, nhưng có giới hạn bão hòa để tránh việc lặp từ khóa quá mức.
        \item \textbf{IDF (Inverse Document Frequency):} Độ hiếm của từ khóa trên toàn bộ tập dữ liệu. Các từ khóa xuất hiện trong ít tài liệu (như thuật ngữ chuyên ngành) sẽ có trọng số cao hơn so với các từ thông dụng (như "the", "a", "is").
    \end{itemize}
\end{itemize}

\section{Field Types trong Elasticsearch}
Elasticsearch phân biệt rõ ràng hai loại dữ liệu văn bản phục vụ cho các mục đích truy vấn khác nhau \cite{elasticsearch_docs}. Bảng \ref{tab:text-vs-keyword} tóm tắt sự khác biệt:

\begin{table}[H]
\centering
\begin{tabular}{|p{3cm}|p{5.5cm}|p{5.5cm}|}
\hline
\textbf{Đặc điểm} & \textbf{\texttt{text}} & \textbf{\texttt{keyword}} \\ \hline
\textbf{Phân tích cú pháp} & Có (tokenization, lowercasing, stop-word removal) & Không -- lưu nguyên bản \\ \hline
\textbf{Cấu trúc lưu trữ} & Inverted Index (ánh xạ term $\to$ documents) & Doc Values (ánh xạ document $\to$ giá trị gốc) \\ \hline
\textbf{Loại truy vấn} & Full-text search (BM25 scoring) & Exact-match, filtering, aggregation \\ \hline
\textbf{Ví dụ trường} & \texttt{title}, \texttt{abstract}, \texttt{authors} & \texttt{year}, \texttt{categories}, \texttt{id} \\ \hline
\textbf{Ví dụ truy vấn} & ``machine learning'' $\to$ khớp ``Machine Learning for...'' & ``2023'' $\to$ chỉ khớp chính xác ``2023'' \\ \hline
\end{tabular}
\caption{So sánh \texttt{text} và \texttt{keyword} field types trong Elasticsearch}
\label{tab:text-vs-keyword}
\end{table}

Việc lựa chọn đúng field type cho từng trường dữ liệu là yếu tố then chốt ảnh hưởng trực tiếp đến hiệu năng tìm kiếm và tính chính xác của kết quả trả về.

\section{Vector Embedding và k-Nearest Neighbors (kNN)}
Để giải quyết hạn chế của tìm kiếm từ khóa (không hiểu ngữ nghĩa), hệ thống đề xuất sử dụng các mô hình ngôn ngữ học sâu (như \texttt{SentenceTransformers} \cite{reimers2019sentencebert}) để chuyển đổi văn bản thành các vector số thực (dense vectors). Thuật toán k-Nearest Neighbors (kNN) sau đó được áp dụng để tính toán độ tương đồng (ví dụ: Cosine Similarity) giữa vector của câu truy vấn và vector của tài liệu, từ đó tìm ra các kết quả có ngữ nghĩa gần nhất.

\section{Reciprocal Rank Fusion (RRF)}
RRF là một thuật toán xếp hạng (ranking algorithm) kết hợp kết quả từ nhiều phương pháp tìm kiếm khác nhau (ở đây là BM25 và kNN) mà không cần phải chuẩn hóa điểm số (score) của chúng \cite{cormack2009rrf}. Công thức RRF tính điểm cho mỗi document $d$ dựa trên vị trí xếp hạng trong từng hệ thống:

\[
\text{RRF}(d) = \sum_{r \in R} \frac{1}{k + r(d)}
\]

trong đó $R$ là tập các bảng xếp hạng (từ BM25 và kNN), $r(d)$ là vị trí của document $d$ trong bảng xếp hạng $r$, và $k$ là hằng số điều chỉnh (thường $k = 60$). Công thức này giúp đẩy các tài liệu xuất hiện ở thứ hạng cao trong cả hai phương pháp lên vị trí top đầu, mang lại kết quả tối ưu và chính xác nhất cho hệ thống Hybrid Search.